% ============================================================
% Title supplied by appendix wrapper.
\label{sec:boundary-defect}
% ============================================================

This section formalizes the ``kernel defect'' as a projection induced by
constraint maps. It supplies the finite-dimensional projector calculus used in
Section~\ref{sec:observable-relative-closure} to split a returned residual into
its workload-effective and workload-null parts. Everything below is stated for
finite-dimensional inner-product spaces, so the Moore--Penrose pseudoinverse
and trace are unambiguous.

\subsection{Constraint Maps}

\begin{definition}[Constraint map]
\label{def:constraint-map}\lean{CoherenceCalculus.AppendixOperators.ConstraintMap}
Let $\cH_{\mathrm{source}}$ and $\cH_{\mathrm{target}}$ be
finite-dimensional inner-product spaces. A \emph{constraint map} is a linear
map
\[
C:\cH_{\mathrm{source}}\to\cH_{\mathrm{target}}
\]
encoding conditions that must be satisfied. The kernel $\Ker(C)$ consists of
states that violate no constraints; the cokernel measures redundancy or defect.
\end{definition}

\subsection{Defect Projectors}

\begin{definition}[Defect projector]
\label{def:defect-projector}\lean{CoherenceCalculus.AppendixOperators.defectProjector}
Define the \emph{defect projector}
\[
P_{\Ker} := I - C^+C,
\]
where $C^+$ is the Moore--Penrose pseudoinverse of $C$.
\end{definition}

\begin{lemma}[Trace equals kernel dimension]
\label{lem:trace-ker}\lean{CoherenceCalculus.AppendixOperators.defect_projector_characterization}
$P_{\Ker}$ is the orthogonal projector onto $\Ker(C)$ and
\[
\Tr(P_{\Ker}) = \dim\Ker(C).
\]
\end{lemma}

\begin{proof}
In finite-dimensional inner-product spaces, $C^+C$ is the orthogonal projector
onto $\Img(C^\ast) = (\Ker C)^\perp$. Therefore
\[
P_{\Ker}=I-C^+C
\]
is the orthogonal projector onto $\Ker(C)$. The trace of an orthogonal
projector equals the dimension of its range, so
\[
\Tr(P_{\Ker})=\dim\Ker(C).
\]
\end{proof}

\subsection{Interpretation and Extensions}

\begin{quote}\small
\textbf{Intuition.} The defect projector $P_{\Ker}$ picks out exactly those
states that have no image under the constraint map $C$. If $\dim\Ker(C) = 1$,
there is a single unmatched degree of freedom, namely one mode that cannot be
matched to any target.
\end{quote}

\begin{remark}[Interpretation: unmatched modes]
\lean{CoherenceCalculus.AppendixOperators.boundary_homeless_modes}
When $\dim\Ker(C) > 0$, there exist source configurations with no valid target representation.
These unmatched modes constitute the defect. A rank-one kernel corresponds to
a single unclosable degree of freedom.
\end{remark}

\begin{remark}[Infinite-dimensional extension]
\lean{CoherenceCalculus.AppendixOperators.boundary_infinite_dimensional_extension}
For closed-range operators between Hilbert spaces the same formulas hold with
the Hilbert-space Moore--Penrose pseudoinverse. That functional-analytic
extension is standard, but the current appendix only uses the exact
finite-dimensional version above.
\end{remark}

\begin{remark}[Boundary defect functional]
\lean{CoherenceCalculus.AppendixOperators.boundary_defect_functional}
In graph or geometric settings where constraint maps are localized (e.g., at boundary events), one can sum the local kernel dimensions:
\[
\cB := \sum_{v \in S} \dim\Ker(C_v).
\]
This ``boundary defect functional'' counts total unclosable degrees of freedom.
\end{remark}
